\chapter{University-Style Mock Examinations}
\label{chap:mock-examinations}

\section{Examination Strategy \& Time Management}
University end-term examinations for ECE249 are typically structured for **100 Marks** over a **3-Hour (180 Minutes)** duration.

\begin{examtip}[Optimal 180-Minute Exam Time Allocation]
\begin{itemize}
    \item \textbf{First 5 Minutes:} Question paper scan and selection of optimal optional questions.
    \item \textbf{Section A (20 Marks):} 10 Short questions $\times$ 2.5 min = \textbf{25 Minutes}.
    \item \textbf{Section B (30 Marks):} 6 Analytical problems $\times$ 9 min = \textbf{54 Minutes}.
    \item \textbf{Section C (50 Marks):} 5 Comprehensive long questions $\times$ 18 min = \textbf{90 Minutes}.
    \item \textbf{Final 10 Minutes:} Numerical verification, unit checking, circuit label inspection.
\end{itemize}
\end{examtip}

---

\section{Mock Examination Paper 1}
\noindent\textbf{Course Code:} ECE249 \hfill \textbf{Maximum Marks:} 100 \\
\textbf{Course Title:} Digital Electronics \& Microcontroller Interfacing \hfill \textbf{Time Allowed:} 3 Hours

\vspace{0.4cm}
\noindent\rule{\textwidth}{1.2pt}

\subsection*{SECTION A: Short Conceptual Questions ($10 \times 2 = 20$ Marks)}
\noindent\textit{Answer all questions. Each question carries 2 marks.}

\begin{enumerate}[label=\textbf{Q\arabic*.}, leftmargin=2em]
    \item State Ohm's Law and name two non-ohmic devices where it fails.
    \item State Kirchhoff's Current Law (KCL) and identify the fundamental physical conservation law upon which it is based.
    \item Distinguish between Intrinsic and Extrinsic semiconductors with respect to charge carrier concentration.
    \item Define the knee voltage of a PN junction diode. State its typical values for Silicon and Germanium.
    \item Why does the Arduino UNO operate with an on-board crystal oscillator frequency of $\SI{16}{\mega\hertz}$?
    \item Convert the binary number $(101101)_2$ into its corresponding Gray code.
    \item State De Morgan's first and second theorems algebraically.
    \item Why are NAND and NOR gates designated as ``Universal Gates''?
    \item Define the Race-Around Condition in a level-triggered JK flip-flop.
    \item Differentiate between a Ring Counter and a Johnson Counter in terms of their modulus for an $N$-flip-flop register.
\end{enumerate}

\subsection*{SECTION B: Analytical \& Numerical Problems ($6 \times 5 = 30$ Marks)}
\noindent\textit{Answer any six questions. Each question carries 5 marks.}

\begin{enumerate}[label=\textbf{Q\arabic*.}, leftmargin=2em, start=11]
    \item Use the Voltage Division and Current Division rules to find the voltage $V_x$ and current $I_x$ in the circuit containing a $\SI{48}{\volt}$ source with series resistors $\SI{12}{\ohm}$ and a parallel combination of $\SI{20}{\ohm}$ and $\SI{30}{\ohm}$.
    \item A full-wave bridge rectifier with ideal diodes delivers power to a load resistance $R_L = \SI{1}{\kilo\ohm}$ from an AC source $v_s(t) = 120\sqrt{2}\sin(120\pi t)\,\si{\volt}$. Calculate: (a) $V_{dc}$, (b) $I_{dc}$, (c) Ripple voltage $V_{ac}$, and (d) Rectifier efficiency $\eta$.
    \item An Arduino UNO ADC reads an analog voltage from an LDR connected in a pull-down voltage divider with a $\SI{10}{\kilo\ohm}$ fixed resistor. If the ADC registers a digital value of $750$ with $V_{ref} = \SI{5.0}{\volt}$, calculate the measured voltage and the instantaneous resistance of the LDR.
    \item Perform the subtraction $(42)_{10} - (68)_{10}$ using 8-bit 2's complement arithmetic. Clearly state the carry status and final signed magnitude.
    \item Minimize the 4-variable logic function using a K-Map:
    $$F(A,B,C,D) = \sum m(1, 3, 7, 11, 15) + d(0, 2, 5)$$
    \item Synthesize a Half Subtractor using only two-input NAND gates.
    \item Convert an available D Flip-Flop into a T Flip-Flop. Show the complete conversion table, K-Map, and circuit diagram.
\end{enumerate}

\subsection*{SECTION C: Comprehensive Design Questions ($5 \times 10 = 50$ Marks)}
\noindent\textit{Answer all five questions. Each question carries 10 marks.}

\begin{enumerate}[label=\textbf{Q\arabic*.}, leftmargin=2em, start=18]
    \item \textbf{(a)} Explain the construction and operation of an NPN Bipolar Junction Transistor in Common Emitter (CE) configuration. \\
    \textbf{(b)} Sketch and explain the input and output V-I characteristics, clearly demarcating the Cut-off, Active, and Saturation regions. \\
    \textbf{(c)} Derive the mathematical relationship between $\alpha$ and $\beta$.
    
    \item \textbf{(a)} Draw the architectural block diagram of an Ultrasonic Distance Sensor (HC-SR04) and explain its operational timing protocol. \\
    \textbf{(b)} Derive the exact distance formula and explain why division by 2 is required. \\
    \textbf{(c)} If the echo pulse width measured by an Arduino microcontroller is $\SI{2500}{\micro\second}$, calculate the physical obstacle distance in centimeters (assuming $v = \SI{340}{\meter\per\second}$).
    
    \item \textbf{(a)} Design a 4:1 Multiplexer using basic logic gates and write its complete Boolean equation. \\
    \textbf{(b)} Implement the 4-variable Boolean function $F(A,B,C,D) = \sum m(0, 1, 3, 4, 8, 9, 15)$ using an 8:1 Multiplexer with $A, B, C$ as select lines.
    
    \item \textbf{(a)} What is the Race-Around condition in JK flip-flops? Derive the mathematical inequality that causes it. \\
    \textbf{(b)} Explain with a complete schematic and timing diagrams how a Master-Slave JK Flip-Flop permanently eliminates the race-around condition.
    
    \item Design a 3-bit Synchronous UP Counter using JK Flip-Flops. \\
    \textbf{(a)} Draw the State Diagram and State Transition Table. \\
    \textbf{(b)} Map the flip-flop excitations using the JK excitation table. \\
    \textbf{(c)} Simplify all $J$ and $K$ input equations using K-Maps. \\
    \textbf{(d)} Draw the final logic circuit schematic.
\end{enumerate}

\vspace{0.4cm}
\noindent\rule{\textwidth}{1.2pt}

---

\section{Solutions for Mock Examination Paper 1}

\subsection*{Solutions for Section A}
\begin{solutionbox}[Solutions to Questions 1–10]
\begin{enumerate}[label=\textbf{Ans Q\arabic*.}]
    \item \textbf{Ohm's Law:} At constant temperature, current through a conductor is proportional to applied voltage ($V = IR$). Two non-ohmic devices: **PN Junction Diode** and **Bipolar Junction Transistor (BJT)**.
    \item \textbf{KCL:} $\sum I_{\text{entering}} = \sum I_{\text{leaving}}$ at any node. It is based on the **Conservation of Electric Charge**.
    \item \textbf{Semiconductors:} In Intrinsic semiconductors, $n = p = n_i$. In Extrinsic N-type, electrons dominate ($n \approx N_D \gg p$); in P-type, holes dominate ($p \approx N_A \gg n$).
    \item \textbf{Knee Voltage ($V_k$):} The forward bias voltage at which diode current begins to increase exponentially. $\mathbf{V_k(\text{Si}) = 0.7\text{ V}}, \mathbf{V_k(\text{Ge}) = 0.3\text{ V}}$.
    \item \textbf{16 MHz Clock:} Established by a quartz crystal oscillator to provide precise instruction cycle execution timing ($\SI{62.5}{\nano\second}$ per clock cycle) for the ATmega328P processor.
    \item \textbf{Binary to Gray:} Binary $101101_2 \implies G_5=1, G_4=1\oplus 0=1, G_3=0\oplus 1=1, G_2=1\oplus 1=0, G_1=1\oplus 0=1, G_0=0\oplus 1=1 \implies \mathbf{111011_{\text{Gray}}}$.
    \item \textbf{De Morgan's Theorems:} (1) $\mathbf{\overline{A \cdot B} = \overline{A} + \overline{B}}$, (2) $\mathbf{\overline{A + B} = \overline{A} \cdot \overline{B}}$.
    \item \textbf{Universal Gates:} Because any Boolean logic function (AND, OR, NOT, XOR, XNOR) can be implemented using exclusively NAND gates or exclusively NOR gates without other gate types.
    \item \textbf{Race-Around Condition:} In a level-triggered JK flip-flop, when $J=K=1$ and clock pulse width $t_p > t_{pd}$ (propagation delay), the output toggles repeatedly and unpredictably during a single clock pulse.
    \item \textbf{Counter Modulus:} An $N$-flip-flop **Ring Counter** has modulus $\mathbf{\text{MOD-}N}$ ($N$ states), whereas an $N$-flip-flop **Johnson Counter** has modulus $\mathbf{\text{MOD-}2N}$ ($2N$ states).
\end{enumerate}
\end{solutionbox}

\subsection*{Solutions for Section B}
\begin{solutionbox}[Solution to Q11: DC Circuit Division]
Parallel equivalent of $\SI{20}{\ohm}$ and $\SI{30}{\ohm}$:
$$R_p = \frac{20 \times 30}{20 + 30} = \frac{600}{50} = \SI{12}{\ohm}$$
Total circuit resistance $R_{total} = R_1 + R_p = 12 + 12 = \SI{24}{\ohm}$.
Total source current $I_s = \frac{V_s}{R_{total}} = \frac{48}{24} = \SI{2.0}{\ampere}$.
By Voltage Division Rule across parallel branch:
$$V_x = V_s \left(\frac{R_p}{R_1 + R_p}\right) = 48 \left(\frac{12}{12 + 12}\right) = \mathbf{\SI{24}{\volt}}$$
By Current Division Rule through the $\SI{20}{\ohm}$ resistor:
$$I_x = I_s \left(\frac{30}{20 + 30}\right) = 2.0 \times \frac{30}{50} = \mathbf{\SI{1.2}{\ampere}}$$
\end{solutionbox}

\begin{solutionbox}[Solution to Q12: Bridge Rectifier Analysis]
Peak secondary voltage $V_m = 120\sqrt{2} \approx \SI{169.7}{\volt}$.
\begin{align*}
    \text{(a) } V_{dc} &= \frac{2V_m}{\pi} = \frac{2 \times 169.7}{\pi} \approx \mathbf{\SI{108.04}{\volt}} \\
    \text{(b) } I_{dc} &= \frac{V_{dc}}{R_L} = \frac{108.04}{\SI{1000}{\ohm}} = \mathbf{\SI{0.108}{\ampere} = \SI{108.04}{\milli\ampere}} \\
    \text{(c) } V_{rms} &= \frac{V_m}{\sqrt{2}} = \SI{120}{\volt} \implies V_{ac} = \sqrt{V_{rms}^2 - V_{dc}^2} = \sqrt{120^2 - 108.04^2} = \mathbf{\SI{52.22}{\volt}} \\
    \text{(d) } \eta &= \frac{8}{\pi^2} \approx \mathbf{81.2\%}
\end{align*}
\end{solutionbox}

\begin{solutionbox}[Solution to Q14: 2's Complement Subtraction]
Compute $(42)_{10} - (68)_{10}$ in 8 bits:
\begin{itemize}
    \item $(42)_{10} = 00101010_2$, $(68)_{10} = 01000100_2$
    \item 1's complement of $68 = 10111011_2 \implies$ 2's complement of $68 = 10111100_2$
    \item Add: $00101010_2 + 10111100_2 = 11100110_2$
    \item No end-around carry ($C_{out} = 0$) $\implies$ Result is Negative.
    \item Magnitude $= 2\text{'s comp of } 11100110_2 = 00011001_2 + 1 = 00011010_2 = (26)_{10}$.
    \item $\mathbf{\text{Result} = -(26)_{10} \quad (11100110_2)}$.
\end{itemize}
\end{solutionbox}

\begin{solutionbox}[Solution to Q15: K-Map Minimization]
$F(A,B,C,D) = \sum m(1, 3, 7, 11, 15) + d(0, 2, 5)$.
\begin{itemize}
    \item Column $CD=11$ contains $m_3, m_7, m_{15}, m_{11} \implies$ forms an **Octet / Quad** $CD$.
    \item Row $AB=00$ contains $d_0, m_1, m_3, d_2 \implies$ forms a **Quad** $\overline{A}\,\overline{B}$.
    \item Remaining $m_5$ (with $m_1, m_3, m_7$) forms Quad $\overline{A}D$.
    \item Minimal equation: $\mathbf{F = \overline{A}\,\overline{B} + CD}$ (or $\overline{A}D + CD$).
\end{itemize}
\end{solutionbox}

\begin{solutionbox}[Solution to Q17: D to T Flip-Flop Conversion]
Target is T-FF ($T, Q_n \rightarrow Q_{n+1}$); Available is D-FF ($Q_{n+1} = D$).
\begin{center}
\small
\begin{tabular}{cc|c|c}
\toprule
$T$ & $Q_n$ & $Q_{n+1}$ & $D$ \\
\midrule
0 & 0 & 0 & 0 \\
0 & 1 & 1 & 1 \\
1 & 0 & 1 & 1 \\
1 & 1 & 0 & 0 \\
\bottomrule
\end{tabular}
\end{center}
From the table, $D = 1$ when $(T, Q_n) = (0,1)$ or $(1,0) \implies \mathbf{D = T \oplus Q_n}$.
Realization: Connect an XOR gate with inputs $T$ and $Q_n$ to the $D$ input terminal.
\end{solutionbox}

---

\section{Mock Examination Paper 2}
\noindent\textbf{Course Code:} ECE249 \hfill \textbf{Maximum Marks:} 100 \\
\textbf{Course Title:} Digital Electronics \& Microcontroller Interfacing \hfill \textbf{Time Allowed:} 3 Hours

\vspace{0.4cm}
\noindent\rule{\textwidth}{1.2pt}

\subsection*{SECTION A: Short Conceptual Questions ($10 \times 2 = 20$ Marks)}
\begin{enumerate}[label=\textbf{Q\arabic*.}, leftmargin=2em]
    \item Why does a silicon PN diode require $\SI{0.7}{\volt}$ to turn ON, whereas Germanium requires only $\SI{0.3}{\volt}$?
    \item State the relationship between collector current $I_C$, base current $I_B$, and current gain $\beta$ in a BJT operating in the active region.
    \item What is the function of the `IOREF` and `AREF` pins on an Arduino UNO board?
    \item Explain the photoconductive principle of a Cadmium Sulfide (CdS) LDR.
    \item Express decimal $29_{10}$ in (a) Pure Binary, and (b) BCD (8421).
    \item Why is Gray code called a ``reflected'' code?
    \item Write the dual of the Boolean algebraic identity: $A + \overline{A}B = A + B$.
    \item Draw the logic symbol and write the Boolean output equation for a 2-input XNOR gate.
    \item What is the primary operational difference between a latch and an edge-triggered flip-flop?
    \item Explain how an asynchronous counter can be truncated to count modulo-6.
\end{enumerate}

\subsection*{SECTION B: Analytical \& Numerical Problems ($6 \times 5 = 30$ Marks)}
\begin{enumerate}[label=\textbf{Q\arabic*.}, leftmargin=2em, start=11]
    \item In a common-emitter NPN transistor circuit, $V_{CC} = \SI{15}{\volt}, R_C = \SI{2.2}{\kilo\ohm}, \beta = 100$. If $I_B = \SI{50}{\micro\ampere}$, calculate $I_C, I_E$ and $V_{CE}$. Verify whether the transistor is in active or saturation mode.
    \item An Arduino UNO ADC pin receives a signal from an ultrasonic sensor. If the measured voltage is $\SI{2.45}{\volt}$, calculate the 10-bit integer returned by `analogRead()`.
    \item Convert Gray code $101101_{\text{Gray}}$ into its binary equivalent.
    \item Using Boolean algebra laws, simplify the expression: $F = A\overline{B}C + \overline{A}\,\overline{B}C + ABC + AB\overline{C}$.
    \item Design a 1-bit binary Half Adder using only 2-input NOR gates.
    \item Implement a Full Subtractor using a 3-to-8 line decoder and external logic gates.
    \item Convert a JK Flip-Flop into an SR Flip-Flop. Show the conversion table and K-Maps.
\end{enumerate}

\subsection*{SECTION C: Comprehensive Design Questions ($5 \times 10 = 50$ Marks)}
\begin{enumerate}[label=\textbf{Q\arabic*.}, leftmargin=2em, start=18]
    \item \textbf{(a)} Explain the working of a Full-Wave Center-Tapped Rectifier with a complete circuit diagram and waveforms. \\
    \textbf{(b)} Derive expressions for its DC output voltage, RMS voltage, ripple factor, and rectification efficiency. \\
    \textbf{(c)} Why is the PIV rating of diodes in a center-tapped rectifier twice that of a bridge rectifier?
    
    \item \textbf{(a)} Detail the pin configuration of the Arduino UNO board, grouping pins by power, analog, digital I/O, PWM, and communications. \\
    \textbf{(b)} Explain how PWM approximates analog output voltages on pins ~3, 5, 6, 9, 10, 11. Calculate the average output voltage for an `analogWrite(6, 153)` command.
    
    \item \textbf{(a)} Minimize the Boolean function using a 4-variable K-Map:
    $$F(A,B,C,D) = \prod M(0, 1, 4, 5, 8, 9, 12, 13) \cdot d(3, 11)$$
    \textbf{(b)} Realize the minimal POS expression using only NOR gates.
    
    \item \textbf{(a)} Explain the architecture and operation of a 4-bit Serial-In Parallel-Out (SIPO) shift register. Draw its schematic and timing diagram. \\
    \textbf{(b)} Explain the operational difference between SISO, SIPO, PISO, and PIPO shift registers.
    
    \item Design a Synchronous 3-bit DOWN Counter ($7 \rightarrow 6 \rightarrow 5 \rightarrow 4 \rightarrow 3 \rightarrow 2 \rightarrow 1 \rightarrow 0 \rightarrow 7$) using T Flip-Flops. \\
    \textbf{(a)} Formulate the State Table and determine $T_2, T_1, T_0$ excitation inputs. \\
    \textbf{(b)} Minimize input expressions using K-Maps. \\
    \textbf{(c)} Draw the complete logic circuit diagram.
\end{enumerate}

\vspace{0.4cm}
\noindent\rule{\textwidth}{1.2pt}

---

\section{Solutions for Mock Examination Paper 2}

\begin{solutionbox}[Selected Solutions to Paper 2]
\textbf{Q11 Solution (BJT Bias):}
\begin{align*}
    I_C &= \beta I_B = 100 \times \SI{50}{\micro\ampere} = \SI{5.0}{\milli\ampere} \\
    I_E &= I_B + I_C = 0.05 + 5.0 = \SI{5.05}{\milli\ampere} \\
    V_{CE} &= V_{CC} - I_C R_C = 15 - (5.0\text{ mA} \times 2.2\text{ k}\Omega) = 15 - 11.0 = \mathbf{\SI{4.0}{\volt}}
\end{align*}
Since $V_{CE} = \SI{4.0}{\volt} > V_{CE(sat)} \approx \SI{0.2}{\volt}$, the transistor is in **Active Mode**.

\vspace{6pt}
\textbf{Q13 Solution (Gray to Binary):}
Gray $= 101101_2$:
\begin{itemize}
    \item $B_5 = G_5 = 1$
    \item $B_4 = B_5 \oplus G_4 = 1 \oplus 0 = 1$
    \item $B_3 = B_4 \oplus G_3 = 1 \oplus 1 = 0$
    \item $B_2 = B_3 \oplus G_2 = 0 \oplus 1 = 1$
    \item $B_1 = B_2 \oplus G_1 = 1 \oplus 0 = 1$
    \item $B_0 = B_1 \oplus G_0 = 1 \oplus 1 = 0$
\end{itemize}
$\mathbf{\text{Binary} = 110110_2}$.

\vspace{6pt}
\textbf{Q19 (b) Solution (PWM Average Voltage):}
$$\text{Duty Cycle} = \frac{153}{255} = 0.60 \implies V_{avg} = 5.0\text{ V} \times 0.60 = \mathbf{\SI{3.0}{\volt}}$$

\vspace{6pt}
\textbf{Q22 Solution (3-Bit Synchronous DOWN Counter):}
Toggle logic for down counter: A flip-flop toggles when all lower-order bits are $0$ ($\overline{Q} = 1$):
$$\mathbf{T_0 = 1, \qquad T_1 = \overline{Q}_0, \qquad T_2 = \overline{Q}_0 \cdot \overline{Q}_1}$$
\end{solutionbox}
